% RESEARCH_DRAFT_SHA256: 4efd15487cfa3c0bc1e31dc08bd43ace3703f805f6f2a6c44bbaa6066dbe71ab
% Layout and source-reference wrappers expanded for the web renderer.
\documentclass[11pt]{article}
\usepackage[margin=25mm]{geometry}
\usepackage{amsthm}
\usepackage{fontspec,unicode-math}
\setmainfont{Latin Modern Roman}
\setmathfont{Latin Modern Math}
\usepackage{amsmath,amssymb,mathtools,xurl,hyperref}
\hypersetup{colorlinks=true,urlcolor=blue,linkcolor=blue}
\setcounter{secnumdepth}{0}
\setlength{\parindent}{0pt}
\setlength{\parskip}{5pt}
\setlength{\emergencystretch}{3em}
\newcommand{\sref}[2]{\hyperlink{source:#1:#2}{[S#2]}}
\newcommand{\cataloguescope}[1]{\paragraph{Source record.}#1}
\newcommand{\floor}[1]{\left\lfloor#1\right\rfloor}
\newtheorem{proposition}{Proposition}
\newtheorem{lemma}[proposition]{Lemma}
\begin{document}
\paragraph{Research manuscript.} OpenAI \texttt{gpt-6-astra}, reasoning \texttt{ultra}; 12 October 2026.
\paragraph{Preserved source body.} The complete catalogue mathematical text and sources follow with layout and source links expanded, including source qualifications and its historical preparation notes. New research begins at the section headed Research attempt.

\section{Duplicate Distribution-Verification Question}\label{statement}

\subsection{Definitions and mathematical statement}

\paragraph{Shared notation.}Standard number systems are $\mathbb N,\mathbb Z,\mathbb Q,\mathbb R,\mathbb C$. Specialized notation and index ranges follow the cited source and the definitions below.


An interactive proof exchanges messages between a prover and verifier, with completeness for true assertions and soundness against dishonest provers. Here the verifier samples an unknown distribution and checks an approximate property claim.
\paragraph{Record scope and qualifications.}
The catalogue title is retained as metadata. Named certificates, contracts, bounds or frameworks in that title are not asserted as results of the original source extract.
The supplied registry flags this extraction as unrecoverable. It is displayed as a literal source fragment; the definitions provide context rather than a replacement formal statement.

Source-record limitation: Source flags the statement as unrecoverable: The official report repeats the two motivating questions already preserved in the earlier verifiable-data-science record and then gives the label-invariant construction.

Related catalogue records: \href{https://mehmetmars7.github.io/Erdosproblems-llm-hunter/problem.html?type=open_problems&id=30006181}{30006181}. Separate source IDs are preserved for their provenance.
The source text below is rendered literally where necessary; damaged formulas and composite extracts are not silently repaired into a new mathematical claim.
The supplied catalogue marks this record as unrecoverable. Its extract and terminology are documented, but a complete original problem has not been reconstructed.
\paragraph{Statement recorded in the supplied source.}
Can the untrusted data server prove that its assertions are approximately correct? Can a short efficiently verifiable proof be generated in polynomial time? Verifying Properties of Distributions: Constructions Tal Herman \href{https://doi.org/10.4171/owr/2025/3}{[S2]}
\paragraph{Source status.}
The supplied catalogue labels this source record partially solved; that metadata was checked against the supplied snapshot on 2026-10-11. The label is preserved for provenance and is not an independent verification of present open status or mathematical correctness. The catalogue literature assessment is dated 2026-08-22.
\subsection{Short English statement}
The supplied fragment concerns duplicate Distribution-Verification Question. Its complete mathematical target is not recovered.
\subsection{Sources}
\begin{itemize}
\item[S1] Ulam AI and the UnsolvedMath Contributors, UnsolvedMath: A Curated Collection of Open Mathematics Problems, catalogue entry. Catalogue curation and metadata: CC BY 4.0. The source extract is retained; definitions, English summaries and qualifications are editorial additions.
\url{https://huggingface.co/datasets/ulamai/UnsolvedMath}.
\item[S2] Cryptography (2025), Oberwolfach Reports, DOI 10.4171/owr/2025/3, pp. 141--174
\url{https://doi.org/10.4171/owr/2025/3}.
\item[S3] Herman and Rothblum, distribution-property arguments
\url{https://arxiv.org/abs/2409.06594}.
\item[S4] Herman and Rothblum, doubly-efficient proofs
\url{https://eccc.weizmann.ac.il/report/2023/161/revision/1/download/}.
\item[S5] OWR 2025/3 Cryptography report
\url{https://ems.press/content/serial-article-files/51349}.
\item Source-specific Oberwolfach Reports licences, including CC BY-SA 4.0 where applicable, are retained; see THIRD\_PARTY\_NOTICES.md.
\end{itemize}
\paragraph{Preparation and review.}
The original source and UnsolvedMath contributors retain attribution for the source extract. Codex (GPT-6; exact serving revision not exposed) assisted with definitions, English summaries, scope notes and formatting. No new human mathematical authorship or independently verified solution is asserted. Catalogue identity checked against the complete supplied registry and current repository. Distinct upstream records, including repeated questions and broad or damaged extracts, are retained for complete coverage of the source CS category. A repeated record is not a new independent conjecture. Complete mathematical-equivalence and current-literature review has not been performed. Related canonical identities: 30006181. Preparation and final review follow contribution prompts C and D, with source gaps disclosed. Structural and rendering checks do not establish a new theorem or open status.
\label{end:statement}

\section{Research attempt}
\subsection{Literature, source context and scope}
The full original catalogue body and source qualifications are preserved above. The relevant primary workshop passages were read in the cached publisher report. Historical claims and subsequent attributed results are distinguished from the mathematical work below. A restricted example, standard lemma or known theorem is not counted as new progress on the entire original target. Missing assumptions and repeated or rhetorical extracts are not silently turned into different solved problems.

The catalogue marks this as a duplicate of the distribution-verification question, and the primary talk reports a general construction for label-invariant properties. We retain the duplicate identity and give a separate worked support-size example. The model grants the honest prover a support description and requires commitment to its list before the verifier takes fresh samples; those are explicit extra hypotheses, not silently inferred from the source.
\subsection{A small-support certificate}
Let $D$ be an unknown distribution on $[N]$, $1\le k<N$, and let the property be support size at most $k$. The prover first supplies a list $S\subseteq[N]$ with $|S|\le k$. Reject malformed lists; otherwise draw $q$ fresh independent samples and accept exactly when all lie in $S$.
\begin{proposition}
Completeness is one when the prover supplies the true support. If $D$ has total-variation distance at least $\varepsilon>0$ from every distribution with support at most $k$, any prover is accepted with probability at most $(1-\varepsilon)^q\le e^{-q\varepsilon}$. Thus $q\ge\lceil\log(1/\delta)/\varepsilon\rceil$ suffices for soundness error $\delta\in(0,1)$.
\end{proposition}
\begin{proof}
For any nonempty $S$, the minimum distance from $D$ to a distribution supported on $S$ is $1-D(S)$. The lower bound is the mass outside $S$, and if $D(S)>0$ it is attained by renormalizing $D$ on $S$; direct summation of absolute differences gives that value. If $D(S)=0$, every distribution on $S$ has distance one. Thus any list of at most $k$ points has $D(S)\le1-\varepsilon$ under the soundness premise. Conditional on the committed list, independence of the fresh samples gives acceptance probability $D(S)^q$. Averaging over a randomized list preserves the bound. The exponential inequality follows from $1-u\le e^{-u}$. Completeness is immediate.
\end{proof}
The list costs $O(k\log N)$ bits; verifying membership can be done by sorting and searching. This is a sample saving paid for with communication and prover information, not an all-resource constant-cost protocol.
\subsection{Why a sample-only verifier cannot simply find the support as cheaply}
Take $N=2k$. The NO distribution is uniform on $[2k]$, at distance $1/2$ from the property. For YES, choose a uniformly random $k$-element subset and use its uniform distribution. For $m\le k$, conditional on seeing no repeated sample, the ordered samples in both experiments are uniform over ordered distinct $m$-tuples: symmetry under permutations of $[2k]$ is transitive on those tuples. The collision probabilities are at most $\binom m2/k$ for YES and $\binom m2/(2k)$ for NO, by a union bound over pairs (also after mixing over the YES support).

If two laws have the same conditional law on an event, their total variation is at most the sum of their probabilities of the event's complement: compare each to that common conditional law. Hence the two sample laws have distance at most $3m(m-1)/(4k)\le m(m-1)/k$. Any sample-only test with completeness $2/3$ for every YES distribution and soundness $1/3$ for NO requires a gap at least $1/3$ also under the YES mixture, and therefore $m(m-1)\ge k/3$. If $m>k$, the order-$\sqrt k$ conclusion already holds. Thus direct testing requires $\Omega(\sqrt k)$ samples in this promise family, whereas a committed support list allows a constant number for fixed error.
\section{Verification and remaining gaps}
Exact small-domain enumeration checks the conditional distinct-sample symmetry, mixture total variation and certificate soundness. The honest prover's ability to obtain the true support from samples alone is not established, especially with arbitrarily rare atoms. The construction does not meet every access, sample or communication requirement in the source's general doubly-efficient theorem, and this duplicate is not an independent solved conjecture.

\section{References and source attribution}
Primary workshop report: \url{https://doi.org/10.4171/owr/2025/3}. Relevant passages read: OWR 2025/3, printed page154, Herman construction talk; related page149 general verification.

All retrievals were checked on 11 October 2026. Original question attribution remains in the preserved Sources section. External theorems are attributed and not claimed as new discoveries or fully verified merely by citation.


\paragraph{Integration review.} Codex (GPT-6; exact serving revision not exposed) reviewed the written arguments and expanded layout and source-reference wrappers. Repository metadata and the production web renderer were checked using Python 3.13 and JavaScript with MathJax 3.2.2. This model-assisted review is not independent human proof verification; no PDF compilation or formal proof checking is asserted.

\section{Completion Estimate}
\noindent\textbf{COMPLETION: 0\%}
No substantiated new progress toward the entire original target is claimed. The worked results and failed route are delimited above; a known result, illustrative model, or auxiliary obstruction is not counted as original-target completion.

\section{Final status and attribution}
\textbf{Attempt status: unresolved.} This describes this research attempt, not a new ruling on the source catalogue's status or on external literature claims.
The mathematics and this self-audit were generated by OpenAI \texttt{gpt-6-astra}, reasoning \texttt{ultra}, on 12 October 2026. Codex (GPT-6; exact serving revision not exposed) orchestrated the batch, read and reviewed the written mathematical arguments, independently reran the finite checks, and contributed formatting/repository review. No human mathematical authorship, independent human proof checking, or formal verification is claimed. The historical preparation paragraphs in the reproduced catalogue body concern that body only.

Contribution prompt A was applied to the complete definition and supplied source record, followed by the adversarial checks of prompt B. Local reproduction files are \texttt{batch026/checks.py}, \texttt{computations.json}, and the per-ID audit JSON. They were not public links at preparation time. Numerical checks are finite and do not prove unrestricted results. No LaTeX engine was found locally; compilation was not performed by this research agent.

\end{document}
